\subsection{Cartan matrix}\label{sga3:XXI:7.3}

Let \(\cR\) be a root datum.  If \(\Delta\) is a system of simple
roots, the \emph{Cartan matrix} of \(\cR\) relative to \(\Delta\)
is the square matrix indexed by \(\Delta\) and defined by
\[
 a_{\alpha,\beta}=(\alpha^\ast,\beta),
 \qquad \alpha,\beta\in\Delta.
\]
If \(\Delta'\) is another system of simple roots and
\(w\in W(\cR)\) satisfies \(w(\Delta)=\Delta'\), then
\[
 (w(\alpha)^\ast,w(\beta))=(\alpha^\ast,\beta).
\]
Thus the Cartan matrix relative to \(\Delta'\) is obtained from
that relative to \(\Delta\) by the isomorphism
\(\Delta\to\Delta'\) of index sets defined by \(w\).  Consequently,
up to isomorphism of its index set, the Cartan matrix depends only
on \(\cR\).

\begin{proposition}\label{sga3:XXI:7.3.1}
The Cartan matrix has the following properties.
\begin{enumerate}
\item[\SGAInlineRefTarget{sga3:XXI:7.3.1:item:i}(i)] \(a_{\alpha,\alpha}=2\), and
  \(a_{\alpha,\beta}\leq0\) for \(\alpha\neq\beta\).
\item[\SGAInlineRefTarget{sga3:XXI:7.3.1:item:ii}(ii)] \(a_{\alpha,\beta}=0\) implies
  \(a_{\beta,\alpha}=0\).
\item[\SGAInlineRefTarget{sga3:XXI:7.2:item:iii}(iii)] There are strictly positive integers
  \(m_\alpha=\operatorname{long}(\alpha)\) such that the matrix
  \[
   (m_\alpha a_{\alpha,\beta})
  \]
  is symmetric, positive, and nondegenerate.
\item[\SGAInlineRefTarget{sga3:XXI:7.2:item:iv}(iv)] The principal minors of
  \((a_{\alpha,\beta})_{\alpha,\beta\in\Delta}\), namely
  \[
   \det(a_{\alpha,\beta})_{\alpha,\beta\in\Delta'}
   \qquad(\Delta'\subset\Delta),
  \]
  are strictly positive.
\item[\SGAInlineRefTarget{sga3:XXI:7.2:item:v}(v)] One has
  \[
   s_\alpha(\beta)=\beta-a_{\alpha,\beta}\alpha,
   \qquad
   s_\alpha(\beta^\ast)
   =\beta^\ast-a_{\beta,\alpha}\alpha^\ast.
  \]
\end{enumerate}
\end{proposition}

\begin{proof}
Part \SGARefLink{sga3:XXI:7.2:item:v}{(v)} is a definition; \SGARefLink{sga3:XXI:7.3.1:item:i}{(i)} follows from
\ref{sga3:XXI:3.2.11}, \SGARefLink{sga3:XXI:7.3.1:item:ii}{(ii)} from \ref{sga3:XXI:2.2.2}, and
\SGARefLink{sga3:XXI:7.2:item:iii}{(iii)} from \ref{sga3:XXI:7.2.9}.  Part \SGARefLink{sga3:XXI:7.2:item:iv}{(iv)} follows immediately
from \SGARefLink{sga3:XXI:7.2:item:iii}{(iii)} and
\[
 \det(m_\alpha a_{\alpha,\beta})_{\alpha,\beta\in\Delta'}
 =
 \left(\prod_{\alpha\in\Delta'}m_\alpha\right)
 \det(a_{\alpha,\beta})_{\alpha,\beta\in\Delta'}.
\]
\end{proof}

\begin{proposition}\label{sga3:XXI:7.3.2}
Let \(\cR\) and \(\cR'\) be two reduced simply connected
(respectively adjoint) root data, let \(\Delta\) and \(\Delta'\)
be systems of simple roots of \(\cR\) and \(\cR'\), and let
\(u:\Delta\to\Delta'\) be an isomorphism.  Denote the corresponding
Cartan matrices by \((a_{\alpha,\beta})\) and
\((a'_{\alpha',\beta'})\), and suppose that
\[
 a'_{u(\alpha),u(\beta)}=a_{\alpha,\beta}.
\]
Then there is a unique isomorphism from \(\cR\) onto \(\cR'\)
that induces \(u\) on \(\Delta\).
\end{proposition}

\begin{proof}
It is enough to give the proof in the adjoint case.  Then
\(M=\Gamma_0(R)\) and \(M'=\Gamma_0(R')\) are the free abelian
groups generated by \(\Delta\) and \(\Delta'\).  Hence there is a
unique group isomorphism \(M\to M'\) inducing \(u\) on \(\Delta\);
denote it again by \(u\).

We show that \(u(R)\subset R'\).  Every root \(\alpha\) of \(R\)
can be written
\[
 \alpha=s_{\alpha_1}\cdots s_{\alpha_n}(\alpha_{n+1}),
 \qquad \alpha_i\in\Delta.
\]
By the hypothesis on \(u\) and the relations
\ref{sga3:XXI:7.3.1}\,\SGARefLink{sga3:XXI:7.2:item:v}{(v)},
\[
 u(\alpha)=
 s_{u(\alpha_1)}\cdots s_{u(\alpha_n)}
 \bigl(u(\alpha_{n+1})\bigr).
\]
It remains to prove that \({}^tu(R'^\ast)\subset R^\ast\).  This
follows from the fact that the elements of \(M^\ast\)
(respectively \(M'^\ast\)) are determined by their duality with
\(R\) or \(\Delta\) (respectively \(R'\) or \(\Delta'\)), by
\ref{sga3:XXI:1.2.5}.
\end{proof}

\begin{corollary}\label{sga3:XXI:7.3.3}
A reduced simply connected or adjoint root datum is determined up
to isomorphism by its Cartan matrix.
\end{corollary}

\begin{corollary}\label{sga3:XXI:7.3.4}
Let \(\cR\) be a reduced simply connected (respectively adjoint)
root datum and let \(\Delta\) be a system of simple roots.  The
group \(E_\Delta(\cR)\) is identified with the group of
automorphisms of the set \(\Delta\) that preserve the Cartan
matrix.
\end{corollary}

\begin{remark}\label{sga3:XXI:7.3.5}
The existence of a root datum corresponding to a given Cartan
matrix satisfying, for example, \SGARefLink{sga3:XXI:7.3.1:item:i}{(i)}, \SGARefLink{sga3:XXI:7.3.1:item:ii}{(ii)}, and \SGARefLink{sga3:XXI:7.2:item:iv}{(iv)} is not easily
settled directly without using the classification.
\end{remark}
